\documentclass[11pt,a4paper]{article}

\usepackage[T1]{fontenc}
\usepackage[utf8]{inputenc}
\usepackage{lmodern}
\usepackage[spanish,es-noquoting,es-tabla]{babel}
\usepackage{amsmath,amssymb,amsthm}
\usepackage{mathtools}
\usepackage{csquotes}
\usepackage{booktabs}
\usepackage{enumitem}
\usepackage[margin=2.3cm]{geometry}
\usepackage{microtype}
\usepackage{xcolor}
\usepackage[colorlinks=true,linkcolor=blue!55!black]{hyperref}

\newcommand{\vect}[1]{\boldsymbol{#1}}
\newcommand{\matr}[1]{\mathbf{#1}}
\newcommand{\code}[1]{\texttt{#1}}
\DeclareMathOperator{\diag}{diag}

\theoremstyle{definition}
\newtheorem{obs}{Observación}

\title{\textbf{Nuestra $g(\vect\psi)$}\\[2mm]
\large La ecuación que se deriva, frente a la del \emph{writeup}}
\author{Repositorio \code{crb}}
\date{5 de septiembre de 2026}

\begin{document}
\shorthandoff{"}
\maketitle

\begin{abstract}
El \emph{writeup} de referencia escribe la medida como una integral continua
de radiancia $L_h$ sobre el p\'ixel y el tiempo. Aqu\'i se escribe
\textbf{la nuestra}: la funci\'on $g(\vect\psi)$ que implementa
\code{simulate\_transient\_torch} m\'as el fondo $b$. Es una sola f\'ormula,
con las tres no suaves a la vista. Despu\'es, la versi\'on suave (la que se
quiere derivar en cerrado) y por qu\'e $\rho$ y $\varphi$ salen correlacionados
en Fisher: el SPAD \textbf{no} est\'a en el origen.
\end{abstract}

%=======================================================================
\section{Qu\'e escribe el \emph{writeup}, y qu\'e escribimos nosotros}

All\'i la cuenta esperada en un bin es del estilo
%
\[
y_q(\vect\Psi)
=\iint L_h(\vect p,t;\vect\Psi)\,d\vect p\,dt,
\]
%
y las derivadas se meten bajo el integrando. Tres reflexiones difusas
($8\pi^3=(2\pi)^3$), p\'ixel como huella con \'area, pulso continuo.

La nuestra es una \textbf{suma de Riemann} sobre los tri\'angulos de la malla
y los \emph{centros} de p\'ixel (sin factor de \'area del p\'ixel), dos
rebotes ($4\pi^2=(2\pi)^2$), pulso $=$ masa en un solo bin. El \'indice de
medici\'on es el par $(p,j)$: p\'ixel $p$, bin $j$. El par\'ametro es
%
\begin{equation}
\vect\psi=(\rho,\varphi).
\end{equation}
%
Eso es lo que se deriva: $\partial g_{p,j}/\partial\vect\psi$.

%=======================================================================
\section{Geometr\'ia que entra en $g$}

La faceta, ya colocada, es una malla de tri\'angulos $t$ con centro
$\vect c_t(\vect\psi)$, \'area $A_t(\vect\psi)$ y normal
$\vect n_t(\vect\psi)$. El l\'aser est\'a en el origen,
$\vect p_L=\vect 0$, con normal $\vect n_L=(0,0,1)$. El suelo tiene normal
$\vect n_F=(0,0,1)$. El p\'ixel $p$ es un \textbf{punto} en el piso,
%
\begin{equation}
\label{eq:pixel}
\vect p_p
=
\bigl(x_p,\, y_p,\, 0\bigr),
\qquad
x_p
=
-\tfrac{\mathrm{FOV}}{2}+\tfrac{\mathrm{FOV}}{2N}+p_x\tfrac{\mathrm{FOV}}{N},
\quad
y_p
=
-\mathrm{FOV}+\tfrac{\mathrm{FOV}}{2N}+p_y\tfrac{\mathrm{FOV}}{N},
\end{equation}
%
con $\mathrm{FOV}=0.25\,\mathrm{m}$, $N=32$. El centro del FOV es
$(0,-\mathrm{FOV}/2)=(0,-0.125)$: \textbf{todo el SPAD vive en $y<0$},
ning\'un p\'ixel cae en el origen ni sobre la arista $x=0$
($N$ par y el muestreo va a centros de celda).

Distancias y cosenos (antes de recortar):
%
\begin{align}
d_{1,t}(\vect\psi)
&=\lVert\vect p_L-\vect c_t\rVert,
&
d_{2,tp}(\vect\psi)
&=\lVert\vect p_p-\vect c_t\rVert,
\label{eq:d}\\[4pt]
\mu_{1,t}
&=\vect n_t\cdot\frac{\vect p_L-\vect c_t}{d_{1,t}},
&
\mu_{2,tp}
&=\vect n_t\cdot\frac{\vect p_p-\vect c_t}{d_{2,tp}},
\label{eq:mu12}\\[4pt]
\mu_{3,t}
&=\vect n_L\cdot\frac{\vect c_t-\vect p_L}{d_{1,t}},
&
\mu_{4,tp}
&=\vect n_F\cdot\frac{\vect c_t-\vect p_p}{d_{2,tp}}.
\label{eq:mu34}
\end{align}
%
Oclusi\'on: recta en $XY$ entre $\vect c_t$ y $\vect p_p$, corte con $x=0$,
%
\begin{equation}
\label{eq:xint}
m_{tp}=\frac{c_{t,y}-y_p}{c_{t,x}-x_p},
\qquad
x_{\mathrm{int},tp}
=
-\frac{c_{t,y}-m_{tp}\,c_{t,x}}{m_{tp}}.
\end{equation}

%=======================================================================
\section{La $g$ dura (el simulador)}
\label{sec:dura}

Con fondo $b=0.1$, intensidad del l\'aser $I_{\mathrm{laser}}=1000$, $[x]_+=\max(0,x)$:
%
\begin{equation}
\label{eq:g-dura}
\boxed{\;
\begin{aligned}
g_{p,j}(\vect\psi)
&=
b
+
\sum_t
A_t(\vect\psi)\,
I_{\mathrm{laser}}
\frac
{[\mu_{1,t}]_+\,[\mu_{2,tp}]_+\,[\mu_{3,t}]_+\,[\mu_{4,tp}]_+}
{4\pi^2\,d_{1,t}(\vect\psi)^2\,d_{2,tp}(\vect\psi)^2}
\\[6pt]
&\hspace{1.6em}
\cdot
\underbrace{\mathbf{1}\bigl[x_{\mathrm{int},tp}(\vect\psi)>0\bigr]}_{\text{oclusi\'on}}
\cdot
\underbrace{\mathbf{1}\!\left[j=\left\lceil\frac{d_{1,t}+d_{2,tp}}{c\,\Delta t}\right\rceil\right]}_{\text{binning}}
\end{aligned}
\;}
\end{equation}
%
Eso es \emph{una} funci\'on $g(\vect\psi)\in\mathbb R^{N^2 N_{\mathrm{bins}}}$
($32\times32\times95=97\,280$ n\'umeros). Las tres no suaves est\'an
expl\'icitas: $\max$, el indicador de $x_{\mathrm{int}}$ y el $\lceil\cdot\rceil$.
El resto (distancias, $1/d^2$, \'areas) ya es $C^\infty$ donde $d>0$.

Es literalmente el bloque de \code{simulate\_transient\_torch}:
\code{intensity} $\times$ \code{noc} depositado en
\code{arrival\_bin = ceil((d1+d2)/(c*bin\_size))}, m\'as $b$ en
\code{forward\_g\_polar}.

%=======================================================================
\section{La $g$ suave (la que se quiere derivar en cerrado)}
\label{sec:suave}

Misma radiometr\'ia y misma geometr\'ia; solo se sustituyen las tres no suaves:
%
\begin{equation}
\label{eq:g-suave}
\boxed{\;
\begin{aligned}
g_{p,j}(\vect\psi)
&=
b
+
\sum_t
A_t\,
I_{\mathrm{laser}}
\frac
{s_\beta(\mu_{1,t})\,s_\beta(\mu_{2,tp})\,s_\beta(\mu_{3,t})\,s_\beta(\mu_{4,tp})}
{4\pi^2\,d_{1,t}^2\,d_{2,tp}^2}
\\[6pt]
&\hspace{1.6em}
\cdot
\sigma\!\bigl(\kappa\, x_{\mathrm{int},tp}\bigr)
\cdot
S_j\!\bigl(t_{0,tp}(\vect\psi)\bigr)
\end{aligned}
\;}
\end{equation}
%
con
%
\begin{align}
s_\beta(x)
&=\frac1\beta\log\bigl(1+e^{\beta x}\bigr)
&&\text{(softplus $\leftrightarrow\max$)},
\label{eq:softplus}\\[4pt]
\sigma(u)
&=\frac1{1+e^{-u}}
&&\text{(sigmoide $\leftrightarrow$ oclusi\'on)},
\label{eq:sig}\\[4pt]
t_{0,tp}
&=\frac{d_{1,t}+d_{2,tp}}{c},
\label{eq:t0}\\[4pt]
S_j(t_0)
&=\frac12\Big[
\mathrm{erf}\frac{(j+1)\Delta t-t_0}{\sqrt2\,\tau}
-
\mathrm{erf}\frac{j\Delta t-t_0}{\sqrt2\,\tau}
\Big]
&&\text{(pulso en el bin $j$ $\leftrightarrow\lceil\cdot\rceil$)}.
\label{eq:Sj}
\end{align}
%
Anchos atados a la discretizaci\'on del simulador:
$\tau=\Delta t/\sqrt{12}$, $\kappa=N/\mathrm{FOV}$ (penumbra de un p\'ixel),
$\beta\sim50$. Al endurecer $(\tau,1/\kappa,1/\beta)\to0$ se recupera
\eqref{eq:g-dura}.

De \eqref{eq:g-suave} el jacobiano $\partial g/\partial\vect\psi$ es cl\'asico
(autograd o sympy). Con Poisson,
%
\begin{equation}
\matr I(\vect\psi)
=
\sum_{p,j}
\frac1{g_{p,j}}
\frac{\partial g_{p,j}}{\partial\vect\psi}
\frac{\partial g_{p,j}}{\partial\vect\psi}^{\!\top},
\qquad
\matr C=\matr I^{+}.
\end{equation}
%
Eso no cambia al suavizar: cambia $g$ y $J$, no la f\'ormula de Fisher.

%=======================================================================
\section{Por qu\'e $\rho$ y $\varphi$ est\'an correlacionados}
\label{sec:corr}

S\'i: \textbf{que ning\'un p\'ixel est\'e en el origen} (ni en $x=0$) es una
raz\'on central de la correlaci\'on. No es la \'unica, y no tiene que ver con
que $g_{p,j}$ no sea exactamente cero (eso lo evita el fondo $b$).

\subsection{Si el \'unico p\'ixel estuviera en $\vect 0$}

Entonces $d_{1,t}=d_{2,t}=\rho$ (la faceta est\'a a distancia $\rho$ del
origen) y el tiempo de vuelo ser\'ia $t_0=2\rho/c$, \textbf{independiente de
$\varphi$}. El binning no enterar\'ia del \'angulo. $\varphi$ solo entrar\'ia
por la oclusi\'on y los cosenos. El t\'ermino de Fisher
%
\begin{equation}
I_{\rho\varphi}
=
\sum_{p,j}
\frac1{g_{p,j}}
\frac{\partial g_{p,j}}{\partial\rho}
\frac{\partial g_{p,j}}{\partial\varphi}
\end{equation}
%
perder\'ia toda la parte que viene del ToF.

\subsection{Lo que tenemos de verdad}

El SPAD es un parche en $y\in[-0.25,0)$, centrado en $(0,-0.125)$. Con $N=32$
par, los centros \eqref{eq:pixel} \textbf{nunca} caen en $x=0$ ni en
$(0,0)$. Entonces
%
\[
d_{2,tp}
=
\bigl\lVert
(\rho\cos\varphi-x_p,\;\rho\sin\varphi-y_p,\;z_t)
\bigr\rVert
\]
%
depende de \textbf{las dos} coordenadas: girar a $\rho$ fijo acerca o aleja la
faceta del parche (que no est\'a en el origen). El camino $d_1+d_2$ tambi\'en.
Por eso el mismo bin es sensible a $\rho$ y a $\varphi$ a la vez, y
$I_{\rho\varphi}\neq0$.

Adem\'as $x_{\mathrm{int}}$ \eqref{eq:xint} depende de $(c_x,c_y)$, o sea de
$(\rho\cos\varphi,\rho\sin\varphi)$: mover $\rho$ desplaza la sombra, que es el
mecanismo que mide $\varphi$. Esa v\'ia de correlaci\'on existir\'ia aunque el
ToF no acoplase.

El signo que sale en toda la rejilla, $r_{\rho\varphi}<0$, es el de
compensaci\'on: alejar y girar hacia el hueco dejan un transitorio parecido al
primer orden. Eso no lo inventa la FD ni lo quita el suavizado
\eqref{eq:g-suave}: la geometr\'ia es la misma.

\subsection{Qu\'e no es la causa}

\begin{itemize}[leftmargin=*]
  \item \textbf{El fondo $b=0.1$.} Impide $g_q=0$ y regulariza $1/g_q$. No
        crea $I_{\rho\varphi}$. Con $b=0$ (y solo los bins con se\~nal) la
        correlaci\'on seguir\'ia ah\'i.
  \item \textbf{Que $x=0$ no tenga p\'ixel.} Evita el $m=\infty$ de
        \eqref{eq:xint} cuando $N$ es impar y hay columna en la arista. Es
        numerica; la correlaci\'on vendr\'ia igual con un p\'ixel en $x=0$
        (ese p\'ixel, de hecho, estar\'ia justo en la frontera de oclusi\'on).
\end{itemize}

\begin{obs}
$\rho$ y $\varphi$ son coordenadas independientes. Lo correlacionado son los
\emph{errores de estimarlas}, porque los datos (ToF al parche desplazado +
sombra) no las separan. Suavizar \eqref{eq:g-dura} en \eqref{eq:g-suave} no
destruye ese t\'ermino cruzado de Fisher.
\end{obs}

\end{document}
